% ============================================================
%  GUÍA DE ESTUDIO — Días 4 al 6
%  Estudiante : Yandri Alexander Piscocama Jaramillo
%  Plan       : Inglés desde Cero — 6 Meses
%  Mes 1      : Fundamentos: Sonidos, Letras y Primeras Palabras
%  Semana 1   : Números, días de la semana y primera pregunta
%  Generada   : Mayo 2026
% ============================================================

\chapter*{Guía de Estudio — Días 4 al 6}
\addcontentsline{toc}{chapter}{Guía de Estudio — Días 4 al 6}
\label{ch:guia_dia4_6}

\noindent
\begin{tabular}{@{}ll@{}}
  \textbf{Estudiante:} & Yandri Alexander Piscocama Jaramillo \\
  \textbf{Período:}    & Días 4 -- 6 de 180 \\
  \textbf{Mes/Semana:} & Mes 1 · Semana 1 \\
  \textbf{Tema:}       & Números del 4 al 10, los días de la semana
                         y tu primera pregunta en inglés \\
\end{tabular}

\bigskip
\noindent
\textbf{¿Qué vas a aprender estos 3 días?}
Completar los números del 4 al 10, aprender los 7 días de la
semana y formular la pregunta más útil para un aprendiz:
\textit{``What is this?''} (¿Qué es esto?).
Al terminar el Día~6 habrás acumulado 30 palabras reales y serás
capaz de decir en inglés cualquier número del 1 al 10 y cualquier
día de la semana sin dudar.

\bigskip
\noindent\rule{\textwidth}{0.4pt}

% ============================================================
\section*{DÍA 4 — Los números del 4 al 8}
\label{sec:dia4}
% ============================================================

\subsection*{Bloque A · 15 min · Vocabulario del Día 4}

\noindent
Antes de empezar: repasa en voz alta los números del 1 al 3 que
aprendiste. \textit{One, two, three.} ¿Los tienes? Bien. Ahora
continúa la secuencia.

\begin{table}[htbp]
  \centering
  \caption{Vocabulario del Día 4 — Números del 4 al 8}
  \label{tab:vocab_dia4}
  \small
  \begin{tabular}{@{}clllp{4.5cm}@{}}
    \toprule
    \textbf{\#} &
    \textbf{Palabra} &
    \textbf{Se pronuncia} &
    \textbf{Significa} &
    \textbf{Ejemplo real} \\
    \midrule
    16 & \textit{four}  & ``for'' &
        cuatro &
        \textit{``I study four days a week''} =
        Estudio cuatro días a la semana \\
    \addlinespace
    17 & \textit{five}  & ``faiv'' &
        cinco &
        \textit{``I have five minutes''} =
        Tengo cinco minutos \\
    \addlinespace
    18 & \textit{six}   & ``six'' &
        seis &
        \textit{``Six months of English''} =
        Seis meses de inglés \\
    \addlinespace
    19 & \textit{seven} & ``seven'' (casi como en español) &
        siete &
        \textit{``Seven days in a week''} =
        Siete días en una semana \\
    \addlinespace
    20 & \textit{eight} & ``eit''
        (la \textit{gh} es completamente silenciosa) &
        ocho &
        \textit{``I sleep eight hours''} =
        Duermo ocho horas \\
    \bottomrule
  \end{tabular}
\end{table}

\noindent
\textbf{Advertencia de pronunciación:}
La palabra \textit{eight} se escribe con cuatro letras que no se
pronuncian (\textit{-ight}).  Así funciona el inglés: lo que ves
no siempre es lo que dices.  La regla de oro es: aprende la
pronunciación junto con la escritura, nunca por separado.

\subsection*{Bloque B · 15 min · Práctica Activa del Día 4}

\noindent
\textbf{Ejercicio: Completa la secuencia y escribe la palabra en inglés.}

\medskip
\begin{tabular}{@{}llll@{}}
  1 = \textit{one}  & 2 = \textit{two}   &
  3 = \textit{three} & 4 = \_\_\_\_\_ \\[4pt]
  5 = \_\_\_\_\_     & 6 = \_\_\_\_\_    &
  7 = \_\_\_\_\_     & 8 = \_\_\_\_\_ \\
\end{tabular}

\bigskip
\noindent
\textbf{Ejercicio 2: Traduce estas oraciones al español.}

\begin{enumerate}
  \item \textit{I have four brothers and two sisters.}
  \item \textit{The book is good and I have five of them.}
  \item \textit{Six and two is eight.}
  \item \textit{You have seven days.}
  \item \textit{I am not eight years old.}
\end{enumerate}

\medskip
\noindent
\textbf{Respuestas:}
\begin{enumerate}
  \item Tengo cuatro hermanos y dos hermanas.
  \item El libro es bueno y tengo cinco de ellos.
  \item Seis más dos es ocho.
  \item Tienes siete días.
  \item No tengo ocho años.
\end{enumerate}

\subsection*{Bloque C · 10 min · Repaso de los Días 1 al 3}

\noindent
Cubre el español. Di la respuesta en voz alta. Marca las que
recuerdas con $\checkmark$ y las que fallas con $\times$.

\begin{table}[htbp]
  \centering
  \caption{Repaso acumulado — 15 palabras de los Días 1 al 3}
  \label{tab:repaso_dia4}
  \small
  \begin{tabular}{@{}llll@{}}
    \toprule
    \textbf{Inglés} & \textbf{Español} &
    \textbf{Inglés} & \textbf{Español} \\
    \midrule
    \textit{the}   & el / la &
    \textit{one}   & uno \\
    \textit{is}    & es / está &
    \textit{two}   & dos \\
    \textit{I}     & yo &
    \textit{three} & tres \\
    \textit{you}   & tú &
    \textit{have}  & tener \\
    \textit{and}   & y &
    \textit{not}   & no (negación) \\
    \textit{hello} & hola &
    \textit{good}  & bueno / bien \\
    \textit{red}   & rojo &
    \textit{blue}  & azul \\
    \textit{yes}   & sí & & \\
    \bottomrule
  \end{tabular}
\end{table}

\subsection*{Bloque D · 5 min · Registro en cuaderno}

\noindent
Fecha + ``Día 4'' + las 5 palabras nuevas + esta oración tuya:
usa un número del 4 al 8 y al menos dos palabras de días
anteriores.  Ejemplo: \textit{``I have six blue books and you
have four red ones.''}

\bigskip
\noindent\rule{\textwidth}{0.4pt}

% ============================================================
\section*{DÍA 5 — Los números 9 y 10, y los primeros días de la semana}
\label{sec:dia5}
% ============================================================

\subsection*{Bloque A · 15 min · Vocabulario del Día 5}

\begin{table}[htbp]
  \centering
  \caption{Vocabulario del Día 5 — Números 9 y 10 + primeros días de la semana}
  \label{tab:vocab_dia5}
  \small
  \begin{tabular}{@{}clllp{4.5cm}@{}}
    \toprule
    \textbf{\#} &
    \textbf{Palabra} &
    \textbf{Se pronuncia} &
    \textbf{Significa} &
    \textbf{Ejemplo real} \\
    \midrule
    21 & \textit{nine}      & ``nain'' &
        nueve &
        \textit{``I study at nine''} =
        Estudio a las nueve \\
    \addlinespace
    22 & \textit{ten}       & ``ten'' &
        diez &
        \textit{``Ten is a good number''} =
        Diez es un buen número \\
    \addlinespace
    23 & \textit{Monday}    & ``mandei'' &
        lunes &
        \textit{``Monday is the first day''} =
        El lunes es el primer día \\
    \addlinespace
    24 & \textit{Tuesday}   & ``tyusdei'' &
        martes &
        \textit{``I study on Tuesday''} =
        Estudio el martes \\
    \addlinespace
    25 & \textit{Wednesday} & ``wensdei''
        (la \textit{d} del medio es silenciosa) &
        miércoles &
        \textit{``Wednesday is day three of the week''} =
        El miércoles es el día tres de la semana \\
    \bottomrule
  \end{tabular}
\end{table}

\noindent
\textbf{Truco de memoria para los días de la semana:}
En inglés, los días terminan casi todos en \textit{-day}
(que suena ``dei'' y significa ``día'').
\textit{Mon-day} = día de la luna, \textit{Tues-day} = día de
Marte (en inglés antiguo), \textit{Wednes-day} = día de Odín.
Conocer el origen te ayuda a recordar la escritura aunque la
pronunciación sea diferente.

\subsection*{Bloque B · 15 min · Práctica Activa del Día 5}

\noindent
\textbf{Ejercicio 1: Escribe el número en inglés.}

\medskip
\begin{tabular}{@{}llll@{}}
  9  = \_\_\_\_\_ &
  10 = \_\_\_\_\_ &
  7  = \_\_\_\_\_ &
  4  = \_\_\_\_\_ \\[4pt]
  3  = \_\_\_\_\_ &
  8  = \_\_\_\_\_ &
  1  = \_\_\_\_\_ &
  6  = \_\_\_\_\_ \\
\end{tabular}

\bigskip
\noindent
\textbf{Ejercicio 2: Completa con el día de la semana correcto.}

\begin{enumerate}
  \item El primer día de la semana laboral es \_\_\_\_\_.
  \item El día que sigue al lunes es \_\_\_\_\_.
  \item \textit{``I study on \_\_\_\_\_ and \_\_\_\_\_''}
        (lunes y miércoles).
  \item \textit{``\_\_\_\_\_ is day two of the week.''}
  \item \textit{``The \_\_\_\_\_ class is good.''}
        (clase del miércoles)
\end{enumerate}

\medskip
\noindent
\textbf{Respuestas:}
\begin{enumerate}
  \item \textit{Monday}
  \item \textit{Tuesday}
  \item \textit{I study on Monday and Wednesday.}
  \item \textit{Tuesday is day two of the week.}
  \item \textit{The Wednesday class is good.}
\end{enumerate}

\subsection*{Bloque C · 10 min · Repaso de los Días 2 al 4}

\noindent
Enfócate especialmente en las palabras que marcaste con $\times$
ayer. Léelas en voz alta tres veces antes de continuar.

\begin{table}[htbp]
  \centering
  \caption{Repaso dirigido — palabras de los Días 2 al 4}
  \label{tab:repaso_dia5}
  \small
  \begin{tabular}{@{}llll@{}}
    \toprule
    \textbf{Inglés} & \textbf{Español} &
    \textbf{Inglés} & \textbf{Español} \\
    \midrule
    \textit{have}  & tener &
    \textit{four}  & cuatro \\
    \textit{not}   & no (negación) &
    \textit{five}  & cinco \\
    \textit{hello} & hola &
    \textit{six}   & seis \\
    \textit{good}  & bueno / bien &
    \textit{seven} & siete \\
    \textit{red}   & rojo &
    \textit{eight} & ocho \\
    \textit{blue}  & azul &
    \textit{yes}   & sí \\
    \textit{one}   & uno &
    \textit{two}   & dos \\
    \textit{three} & tres & & \\
    \bottomrule
  \end{tabular}
\end{table}

\subsection*{Bloque D · 5 min · Registro en cuaderno}

\noindent
Fecha + ``Día 5'' + las 5 palabras nuevas + esta oración:
\textit{``On Monday and Wednesday I study ten words.''}
(o la tuya propia, usando al menos un día de la semana y un número)

\bigskip
\noindent\rule{\textwidth}{0.4pt}

% ============================================================
\section*{DÍA 6 — El resto de la semana y tu primera pregunta}
\label{sec:dia6}
% ============================================================

\subsection*{Bloque A · 15 min · Vocabulario del Día 6}

\begin{table}[htbp]
  \centering
  \caption{Vocabulario del Día 6 — Días restantes y primera pregunta}
  \label{tab:vocab_dia6}
  \small
  \begin{tabular}{@{}clllp{4.5cm}@{}}
    \toprule
    \textbf{\#} &
    \textbf{Palabra} &
    \textbf{Se pronuncia} &
    \textbf{Significa} &
    \textbf{Ejemplo real} \\
    \midrule
    26 & \textit{Thursday} & ``zersdei'' &
        jueves &
        \textit{``Thursday is day four''} =
        El jueves es el día cuatro \\
    \addlinespace
    27 & \textit{Friday}   & ``fraidei'' &
        viernes &
        \textit{``Friday is a good day''} =
        El viernes es un buen día \\
    \addlinespace
    28 & \textit{Saturday} & ``saturdei'' &
        sábado &
        \textit{``I study on Saturday''} =
        Estudio el sábado \\
    \addlinespace
    29 & \textit{Sunday}   & ``sandei'' &
        domingo &
        \textit{``Sunday is day seven''} =
        El domingo es el día siete \\
    \addlinespace
    30 & \textit{what}     & ``wat'' &
        qué / cuál &
        \textit{``What is this?''} =
        ¿Qué es esto? \\
    \bottomrule
  \end{tabular}
\end{table}

\noindent
\textbf{Hito importante:} Con \textit{what} aprendes hoy tu
primera palabra interrogativa en inglés.  Combinada con \textit{is}
(que ya conoces desde el Día~1), puedes construir la pregunta
más útil del idioma: \textit{``What is this?''}
Úsala con cualquier objeto que no conozcas en inglés.

\subsection*{Bloque B · 15 min · Práctica Activa del Día 6}

\noindent
\textbf{Ejercicio 1: Completa la semana completa en orden.}

\medskip
\begin{tabular}{@{}lll@{}}
  Día 1 = \_\_\_\_\_ (lunes) &
  Día 2 = \_\_\_\_\_ (martes) &
  Día 3 = \_\_\_\_\_ (miércoles) \\[4pt]
  Día 4 = \_\_\_\_\_ (jueves) &
  Día 5 = \_\_\_\_\_ (viernes) &
  Día 6 = \_\_\_\_\_ (sábado) \\[4pt]
  Día 7 = \_\_\_\_\_ (domingo) & & \\
\end{tabular}

\bigskip
\noindent
\textbf{Respuestas:}
\textit{Monday, Tuesday, Wednesday, Thursday, Friday, Saturday, Sunday.}

\bigskip
\noindent
\textbf{Ejercicio 2: Escribe en inglés las siguientes oraciones.}

\begin{enumerate}
  \item El viernes es el día cinco de la semana.
  \item El sábado y el domingo son buenos días.
  \item ¿Qué es esto? (usa: \textit{What / is / this / ?})
  \item El jueves no es el día seis.
  \item Estudio el lunes, el miércoles y el viernes.
\end{enumerate}

\medskip
\noindent
\textbf{Respuestas:}
\begin{enumerate}
  \item \textit{Friday is day five of the week.}
  \item \textit{Saturday and Sunday are good days.}
  \item \textit{What is this?}
  \item \textit{Thursday is not day six.}
  \item \textit{I study on Monday, Wednesday, and Friday.}
\end{enumerate}

\subsection*{Bloque C · 10 min · Repaso de los Días 4 y 5}

\noindent
Repasa en voz alta, cubriendo el español. Sin mirar, intenta
decir todos los números del 1 al 10 seguidos:
\textit{one, two, three, four, five, six, seven, eight, nine, ten.}
Si lo haces sin fallar, es tu primer logro completo del plan.

\begin{table}[htbp]
  \centering
  \caption{Repaso acumulado — palabras de los Días 4 y 5}
  \label{tab:repaso_dia6}
  \small
  \begin{tabular}{@{}llll@{}}
    \toprule
    \textbf{Inglés} & \textbf{Español} &
    \textbf{Inglés} & \textbf{Español} \\
    \midrule
    \textit{four}      & cuatro    & \textit{nine}      & nueve \\
    \textit{five}      & cinco     & \textit{ten}       & diez \\
    \textit{six}       & seis      & \textit{Monday}    & lunes \\
    \textit{seven}     & siete     & \textit{Tuesday}   & martes \\
    \textit{eight}     & ocho      & \textit{Wednesday} & miércoles \\
    \bottomrule
  \end{tabular}
\end{table}

\subsection*{Bloque D · 5 min · Registro en cuaderno}

\noindent
Fecha + ``Día 6'' + las 5 palabras nuevas + escribe los 7 días
de la semana de memoria, en orden, sin mirar.
Compara con la respuesta del Ejercicio~1 y encierra en un círculo
cualquier día que hayas escrito mal.

\bigskip
\noindent\rule{\textwidth}{0.4pt}

% ============================================================
\section*{Gramática del Período — Cómo hacer una pregunta básica}
\label{sec:gramatica_periodo2}
% ============================================================

\noindent
\textbf{Concepto:} La pregunta con \textit{What is\ldots?}

\medskip\noindent
En español preguntas ``¿Qué es esto?'' invirtiendo la entonación.
En inglés la estructura es diferente: la palabra interrogativa
va \textbf{siempre al inicio} de la oración.

\medskip
\begin{center}
  \large
  \textbf{What} \quad + \quad \textbf{is} \quad + \quad
  \textbf{esto / aquello / el sujeto}
\end{center}

\medskip\noindent
\textbf{3 ejemplos concretos:}

\begin{enumerate}
  \item \textit{``What is this?''} \quad $\rightarrow$ \quad
        ¿Qué es esto? (señalando algo cercano)
  \item \textit{``What is your name?''} \quad $\rightarrow$ \quad
        ¿Cuál es tu nombre?
  \item \textit{``What is Monday in Spanish?''} \quad $\rightarrow$ \quad
        ¿Qué es \textit{Monday} en español?
\end{enumerate}

\medskip\noindent
\textbf{Error común de hispanohablantes:}
Tendemos a decir \textit{``What is?''} sin completar la oración,
como si fuera suficiente.  En inglés la pregunta siempre necesita
un sujeto o complemento después del verbo.
\textit{``What is?''} $\rightarrow$ incorrecto. \quad
\textit{``What is this?''} $\rightarrow$ correcto.

\medskip\noindent
\textbf{Aplicación práctica inmediata:}
Desde hoy, cuando veas un objeto en tu casa o en la calle cuyo
nombre en inglés no conozcas, dite a ti mismo en voz baja:
\textit{``What is this?''} y búscalo.  Así aprendes vocabulario
nuevo sin listas aburridas.

% ============================================================
\section*{Pronunciación del Período — Vocales largas vs.\ cortas}
\label{sec:pronunciacion_periodo2}
% ============================================================

\noindent
El inglés tiene un sistema de vocales diferente al español.
En español todas las vocales suenan igual siempre.
En inglés la misma letra puede sonar de dos formas distintas:
larga o corta.  Esta es una de las razones por las que el inglés
suena tan diferente a como se escribe.

\medskip\noindent
\textbf{Dos ejemplos con palabras que ya conoces:}

\begin{itemize}
  \item \textit{nine} (nain) $\rightarrow$ la \textit{i} suena larga,
        como la ``ai'' del español.
  \item \textit{six} (six) $\rightarrow$ la \textit{i} suena corta,
        como la ``i'' normal del español.
\end{itemize}

\medskip\noindent
\textbf{Regla práctica (funciona el 70\,\% del tiempo):}
Si una vocal va seguida de consonante + \textit{e} final silenciosa,
la vocal suena larga.  Ejemplo: \textit{five} = ``faiv'' (vocal larga),
\textit{fin} = ``fin'' (vocal corta).

\medskip\noindent
\textbf{Práctica en voz alta:}
Alterna estas dos palabras que ya conoces, sintiendo la diferencia:\\
\textit{nine} (larga) \quad -- \quad \textit{ten} (corta) \quad
-- \quad \textit{five} (larga) \quad -- \quad \textit{six} (corta)

% ============================================================
\section*{Conexión Cognitiva — Lo que acaba de pasar en tu cerebro}
\label{sec:motivacion_periodo2}
% ============================================================

\noindent
En 6 días has acumulado 30 palabras en un idioma que no sabías.
Pero más importante que el número es lo que ocurrió cognitivamente:
tu cerebro construyó un sistema de categorías nuevas, porque aprender
los días de la semana en otro idioma no es simplemente memorizar
etiquetas; es crear una segunda red conceptual de tiempo y calendario
que coexiste con la del español.  Estudios de neuroimagen muestran que
este proceso activa la corteza prefrontal, la misma región responsable
de la toma de decisiones complejas y la planificación a largo plazo.
Cada día que completes el hábito de 45 minutos estás entrenando esa
región, no solo aprendiendo inglés.

% ============================================================
\section*{Vista Previa — Días 7 al 9}
\label{sec:preview_dia7_9}
% ============================================================

\noindent
En la próxima guía aprenderás: los números del 11 al 20,
las 5 reglas de pronunciación más importantes del inglés
(las que te explican por qué las palabras no suenan como se
escriben) y cómo decir la hora básica en inglés.
Con eso terminas la Semana~1 del plan completo.

\bigskip
\begin{center}
  \small\itshape
  Guía generada con el Prompt Reutilizable v1.0 · Mayo 2026 \\
  Siguiente guía: indicar Días 7--9, Mes 1 · Semana 2,
  y adjuntar esta guía como contexto previo.
\end{center}
